\documentclass[12pt,a4paper]{article}
\usepackage{ctex}
\usepackage{geometry}
\geometry{margin=1in}
\usepackage{graphicx}
\usepackage{amsmath,amssymb}
\usepackage{booktabs}
\usepackage{hyperref}
\hypersetup{colorlinks=true, linkcolor=blue, citecolor=blue, urlcolor=blue}
\usepackage{setspace}
\usepackage{fancyhdr}
\pagestyle{fancy}
\lhead{工业软件导学课程笔记}
\rhead{哈尔滨工业大学（深圳）}
\cfoot{\thepage}
% 修复页眉高度警告
\setlength{\headheight}{15pt}

\begin{document}

\begin{titlepage}
    \centering
    \vspace*{1cm}
    \textbf{\Large \color{red} ⚠️ 本文档仅作学习参考使用，严禁抄袭照搬，抄袭后果自负}
    \vspace{2cm}
    \Huge{\textbf{工业软件导学课程笔记}}\\
    \vspace{1cm}
    \Large{课程学习报告}\\
    \vspace{4cm}
    \begin{tabular}{ll}
        \textbf{任课教师}\hspace{2em} & 周晔欣
    \end{tabular}
    \vspace{3cm}
\end{titlepage}

\section*{摘要}
EDA电子设计自动化、CAE计算机辅助工程是电子制造、高端装备研发环节不可或缺的核心工业软件，长期被海外厂商垄断，是制约我国制造业自主可控的关键“卡脖子”领域。本次《工业软件导学》课程邀请三位产业一线专家开展系列专题讲座，分别为嘉立创EDA莫志宏先生、英特仿真行业专家、泊松软件熊老师。三场分享依次从国产开源EDA生态建设、本土自主CAE仿真工程落地、全球车企数字化全生命周期平台架构三个维度，结合企业实战案例、高校产教融合、学科竞赛人才培养等内容展开深度讲解。本文以课堂讲授内容为基础，结合三场专家讲座素材、国内工业软件产业公开白皮书与相关学术文献完成梳理总结，系统梳理EDA、CAE两类工业软件的技术体系、产业现状、国产化痛点与发展路径，同时结合课程学习形成个人认知与实践思考。
\textbf{关键词：}工业软件；EDA；CAE；国产替代；产教融合；数字化研发

\section{课程讲座整体概述}
本次《工业软件导学》课程引入产业专家进课堂的教学形式，围绕国产EDA、CAE两大核心工业软件赛道开展专题分享，打破了传统课堂纯理论的学习模式，从产业一线、工程落地、全球技术对比多个视角补充课堂知识。三场讲座嘉宾研究方向与从业背景高度互补，完整覆盖电子电路设计、多物理场仿真、智能制造数字化平台等高端装备研发全流程工具链，让我系统认识了国内工业软件的发展现状、技术短板与未来突破方向，极大拓宽了我对工业软件以及制造业数字化转型的认知边界。

其中，嘉立创EDA专家讲座聚焦轻量化、开源化国产电子设计工具的普及与人才培养；英特仿真聚焦国产自主CAE仿真软件的工程落地；泊松软件专家结合海外头部车企多年研发平台架构经验，从全球化视角剖析了国内外工业软件体系的差距，三场讲座层层递进、互为补充，形成了完整的产业认知体系。

\section{嘉立创EDA专题讲座学习总结}
\subsection{嘉宾专业背景}
本次EDA专题分享由嘉立创EDA高校教育经理、生态社区负责人莫志宏先生主讲。莫志宏先生为电子信息工程专业科班出身，具备扎实的电子技术专业背景，自2019年起长期负责嘉立创EDA全国高校教育对接、软件培训、赛事指导与开源社区运营管理工作。同时，莫先生深耕全国大学生科创赛事与产教融合领域，担任全国三维数字化创新设计大赛评审专家、蓝桥杯EDA赛项培训专家、iCAN大学生创新创业大赛（黑龙江赛区）评审专家、机械工业出版社专家委员会委员、全国智能物联产教融合共同体副理事长，兼具技术实践能力与高校人才培养经验。

\subsection{讲座核心内容}
莫志宏先生在讲座中重点讲解了国产EDA软件的发展现状、开源生态优势以及面向高校的人才培养体系。在传统电子设计领域，海外商用EDA软件长期垄断市场，诸如Altium Designer、Cadence等主流软件存在授权费用高昂、使用门槛高、本土化适配不足等问题，无论是高校教学实训还是学生自主学习，都存在诸多限制。而嘉立创EDA作为国内领先的开源免费电子设计自动化工具，精准贴合高校教学、学生科创、小微企业轻量化研发场景，实现了本土化功能优化，大幅降低了电子设计的学习与实践门槛。

讲座中提到，嘉立创EDA构建了完善的高校产教融合体系，面向高校提供全套教学课件、实训案例、课程配套资源，助力高校开展电子设计实训教学；面向学生开放永久免费的在线设计工具与海量开源工程案例，为零基础学生入门PCB设计、电路开发提供了便捷平台。同时，平台深度对接蓝桥杯、三维数字化创新设计大赛、iCAN等主流科创赛事，提供专项技术培训、赛事指导与硬件打样支持，打通了“课堂学习—赛事实践—工程落地”的学习链路。

除此之外，嘉立创搭建了国内顶尖的电子设计开源社区，鼓励学生分享开源工程、通用电路模块与设计方案，形成了互助共享的技术学习生态。区别于海外软件以商业盈利为核心的模式，国产EDA软件以人才培养、生态建设为核心，持续培育本土电子设计人才，为我国电子信息产业底层技术人才储备提供了重要支撑。

\section{国产CAE仿真软件行业专题讲座总结}
本次课程开展两场CAE领域专题讲座，分别邀请国内本土CAE领军企业英特仿真专家、泊松软件首席产品专家熊老师授课，从本土软件落地应用、全球行业技术对比两个维度，全方位讲解了国内CAE工业软件的发展现状与技术瓶颈。

\subsection{英特仿真专题分享}
CAE计算机辅助工程软件是高端装备研发的核心工具，可实现力学、热学、流体、电磁等多物理场耦合仿真，广泛应用于电机设计、航空装备、轨道交通、新能源汽车等领域，是制造业数字化研发的核心底座。长期以来，国内高端仿真市场被ANSYS、西门子等海外软件垄断，国产CAE软件起步较晚，存在算法积累不足、多场耦合精度有限、工程案例匮乏等问题。

英特仿真作为国内CAE领域的领军企业，深耕自主仿真求解器研发与本土化工程应用。本次讲座中，企业专家重点介绍了国产自主CAE软件的技术迭代路线、工业落地案例以及校企协同育人模式。相较于海外商用软件，国产CAE软件更贴合国内行业标准与企业定制化需求，本土化服务响应更快、适配性更强，在中端工业仿真场景已经实现规模化替代，同时持续攻坚高端多物理场耦合仿真技术，逐步打破海外技术垄断。

\subsection{泊松软件熊老师专题分享}
\subsubsection{嘉宾从业背景}
熊老师现任泊松软件首席产品管理专家，本科毕业于上海交通大学热能与动力工程专业，后赴德国卡尔斯鲁厄大学攻读车辆工程硕士学位，拥有国内外顶尖工科教育背景。毕业后长期就职于梅赛德斯、西门子戴姆勒等全球头部企业，深耕汽车智能制造与工业软件领域，具备丰富的全球行业架构规划与大型数字化平台落地经验。曾任戴姆勒卡车全球设计与制造流程首席解决方案架构师，主导西门子Teamcenter产品全生命周期数字化管理平台的整体规划、架构设计与全域交付，统筹多团队研发、测试与实施工作，2026年初正式加入泊松软件，赋能国产CAE软件产品升级。

\subsubsection{讲座核心观点}
熊老师结合海外车企成熟的数字化研发体系，深度剖析了国内外工业软件的核心差距。海外高端制造业已形成“CAD建模—CAE仿真—PLM全生命周期管理”一体化数字化闭环，工具链完善、数据互通、流程标准化，研发效率极高。而国内多数制造企业存在设计、仿真、数据管理工具割裂的问题，数据无法高效流转，研发流程碎片化，整体数字化水平仍有较大提升空间。

同时，熊老师提出了国产CAE软件的核心突破思路：国产工业软件的发展不能局限于单一仿真求解器的技术迭代，更需要搭建一体化数字化研发平台，完善产品全生命周期管理体系；依托国内庞大的新能源、高端装备市场，积累海量本土工程数据，持续优化求解算法与仿真精度，通过工程实践反哺技术迭代，逐步缩小与海外成熟软件的差距。

\section{EDA与CAE国产工业软件发展对比与痛点分析}
结合本次课程三场专家讲座内容，可以清晰区分EDA与CAE两类核心工业软件的发展现状与国产化差异。EDA电子设计自动化软件聚焦电路、PCB、芯片设计，国内依托嘉立创等开源工具，在高校教学、中小型企业轻量化设计、科创赛事场景已实现大规模国产替代，生态建设日趋完善，但高端芯片设计EDA工具仍高度依赖海外厂商，核心高端领域仍存在技术壁垒。CAE仿真软件聚焦多物理场数值仿真，技术门槛更高、算法积累周期更长，目前国内企业仍处于追赶阶段，中端场景逐步实现替代，高端多场耦合仿真、大型工业项目仍依赖海外软件。

两类国产工业软件存在共性发展痛点：一是海外软件经过数十年技术迭代，积累了海量算法库、材料库、工程案例库，形成了极高的技术壁垒；二是国内高校传统教学长期依赖海外商用软件，国产软件教学资源、实训体系不够完善，人才培养与国产产业发展脱节；三是工业软件研发投入大、周期长、短期收益低，企业持续研发迭代压力较大。

同时，三场讲座也指明了统一的国产替代路径：以高校产教融合为基础，普及国产工业软件教学与实践，积累本土工程案例与实测数据，培育本土技术人才，通过人才储备与工程实践双向赋能，推动国产工业软件持续迭代升级，逐步实现全领域自主可控。

\section{课程学习收获与个人思考}
通过本次《工业软件导学》课程与三场行业专家讲座学习，我彻底改变了对工业软件的浅层认知。此前我仅将EDA、CAE软件当作课程实训、图纸绘制、仿真计算的工具，并未意识到工业软件是我国高端制造业自主可控的核心底座，是电子、电气、汽车、航空等核心产业的关键基础设施，工业软件的“卡脖子”问题，直接制约着我国高端装备制造业的发展。

本次课程让我清晰认识到国产工业软件的发展现状与机遇：在国家自主可控战略的推动下，以嘉立创EDA、英特仿真、泊松软件为代表的国产企业正在快速崛起，从轻量化、通用化场景逐步突破，持续完善本土软件生态。同时，行业专家的从业经历也让我明白，工业软件的学习不能只停留在软件表层操作，更需要理解底层数学算法、物理原理与工程逻辑，只有掌握核心原理，才能真正适配产业发展需求。

作为电气工程专业的学生，工业软件是未来科研、工程实践、行业就业的核心技能。本次课程学习让我明确了后续学习方向：主动学习国产EDA、CAE工业软件，积极参与国产软件相关科创实践与工程实训，跳出依赖海外软件的固有思维，深耕数字化研发技术，积累本土工业软件实操经验，努力成为适配国内制造业数字化转型、工业软件国产化发展的专业人才。

\section{总结与后续学习规划}
本次系列专题课程通过产业专家的一线分享，让我系统掌握了EDA、CAE工业软件的产业现状、技术特点与国产化发展趋势，深刻理解了工业软件对于国家高端制造产业自主可控的战略意义，弥补了课堂理论与产业实际脱节的短板，建立了完整的工业软件认知体系。

在后续学习中，我将持续深耕工业软件相关知识：第一，强化实操能力，熟练掌握嘉立创EDA等国产工具的电路设计、PCB开发技能，尝试使用国产CAE软件开展基础仿真实践；第二，夯实理论基础，学习有限元仿真、电路设计底层原理，摆脱纯工具化操作的学习局限；第三，积极参与科创赛事与工程实践，依托国产工业软件完成项目开发，积累工程实战经验；第四，持续关注国产工业软件产业动态，跟进国产EDA、CAE软件的技术迭代与生态建设，不断完善自身知识体系，为未来投身智能制造、国产工业软件相关领域打下坚实基础.

\end{document}